% ============================================================
%  HCMUT-DEE Thesis Kit v2.0.3 — Chương 1
%  Copyright (c) 2026 Nguyen Trong Thang & TS. Nguyen Phuc Khai
%  License: MIT
% ============================================================

\chapter{Giới Thiệu Template \& Hướng Dẫn Cài Đặt}
\label{ch:introduction}

% ═══════════════════════════════════════════════════════════════
%  GHI CHU: Chuong nay la HUONG DAN SU DUNG template.
%  Khi viet do an that, hay XOA toan bo noi dung chuong nay
%  va thay bang chuong "Mo dau" cua do an cua ban.
% ═══════════════════════════════════════════════════════════════

\section{Giới thiệu HCMUT-DEE Thesis Kit}

Chào mừng bạn đến với \textbf{HCMUT-DEE Thesis Kit} phiên bản \TemplateVersion{} --- bộ template \LaTeX{} được thiết kế dành riêng cho sinh viên Khoa Điện - Điện Tử, Trường Đại học Bách Khoa TP. Hồ Chí Minh (ĐHQG-HCM).

Template này giúp bạn:
\begin{itemize}
  \item Định dạng đồ án tốt nghiệp \textbf{đúng chuẩn} của trường.
  \item \textbf{Tự động sinh} trang bìa, trang nhiệm vụ, tờ chấm.
  \item Quản lý tài liệu tham khảo chuẩn \textbf{IEEE}.
  \item Chèn mã nguồn Python, MATLAB, C/C++, PLC với \textbf{syntax highlighting}.
  \item Tương thích \textbf{Overleaf}, TeXStudio, VS Code.
\end{itemize}

\section{Cấu trúc thư mục}
\label{sec:structure}

Dưới đây là cấu trúc thư mục của template:

\begin{lstlisting}[style=latexstyle, caption={Cấu trúc thư mục HCMUT-DEE Thesis Kit}, label={lst:tree}]
thesis/
|-- thesis.tex            % File chinh (bien dich file nay)
|-- hcmut-dee.cls         % Class file (KHONG CHINH SUA)
|-- references.bib        % Tai lieu tham khao
|-- front-matter/         % Phan dau (bia, tom tat, ...)
|   |-- acknowledgment.tex
|   |-- abstract-vi.tex
|   |-- abstract-en.tex
|   +-- abbreviations.tex
|-- chapters/             % Noi dung cac chuong
|   |-- chapter1.tex
|   |-- ...
|   +-- chapter6.tex
|-- appendices/           % Phu luc
|   +-- appendix-sample.tex
+-- figures/              % Hinh anh
    |-- logos/             % Logo truong
    +-- examples/         % Hinh mau
\end{lstlisting}

\section{Cài đặt môi trường}
\label{sec:installation}

\subsection{Cách 1: Sử dụng Overleaf (Khuyến nghị)}

Overleaf là nền tảng \LaTeX{} trực tuyến, \textbf{không cần cài đặt} gì trên máy tính.

\begin{enumerate}
  \item Truy cập \url{https://www.overleaf.com} và đăng ký tài khoản (miễn phí).
  \item \textbf{Upload} toàn bộ thư mục \texttt{thesis/} lên Overleaf:
    \begin{itemize}
      \item Chọn ``New Project'' $\rightarrow$ ``Upload Project''.
      \item Nén folder \texttt{thesis/} thành file \texttt{.zip} rồi upload.
    \end{itemize}
  \item Trong Overleaf, vào \textbf{Menu} $\rightarrow$ đặt ``Main document'' là \texttt{thesis.tex}.
  \item Đặt \textbf{Compiler} là \texttt{pdfLaTeX}.
  \item Nhấn \textbf{Recompile} để biên dịch.
\end{enumerate}

\subsection{Cách 2: TeXStudio + MiKTeX (Windows)}

\begin{enumerate}
  \item Cài đặt \textbf{MiKTeX}: \url{https://miktex.org/download}
    \begin{itemize}
      \item Chọn ``Install missing packages on the fly: Yes''.
    \end{itemize}
  \item Cài đặt \textbf{TeXStudio}: \url{https://www.texstudio.org/}
  \item Mở file \texttt{thesis.tex} trong TeXStudio.
  \item Chọn \textbf{Build \& View} (phím tắt: \texttt{F5}).
\end{enumerate}

\subsection{Cách 3: VS Code + LaTeX Workshop}

\begin{enumerate}
  \item Cài đặt TeX Live: \url{https://www.tug.org/texlive/}
  \item Cài extension \textbf{LaTeX Workshop} trong VS Code.
  \item Mở folder \texttt{thesis/} trong VS Code.
  \item Biên dịch bằng \texttt{Ctrl+Alt+B} hoặc lưu file (auto-build).
\end{enumerate}

\section{Biên dịch lần đầu}
\label{sec:first-compile}

Quy trình biên dịch đầy đủ gồm 4 bước:

\begin{enumerate}
  \item \texttt{pdflatex thesis.tex} --- Lần 1: tạo file .aux
  \item \texttt{bibtex thesis} --- Xử lý tài liệu tham khảo
  \item \texttt{pdflatex thesis.tex} --- Lần 2: cập nhật tham chiếu
  \item \texttt{pdflatex thesis.tex} --- Lần 3: hoàn thiện mục lục
\end{enumerate}

\textbf{Lưu ý:} Trên Overleaf, bạn chỉ cần nhấn ``Recompile'' --- hệ thống tự chạy đủ các bước.

\section{Chỉnh sửa thông tin cá nhân}
\label{sec:personal-info}

Mở file \texttt{thesis.tex} và chỉnh sửa \textbf{duy nhất} phần ``THÔNG TIN SINH VIÊN'':

\begin{lstlisting}[style=latexstyle, caption={Các lệnh thông tin cá nhân trong thesis.tex}]
\ThesisTitle{Ten de tai tieng Viet}
\ThesisTitleEN{Thesis Title in English}
\StudentName{Ho Ten Sinh Vien}
\StudentID{21xxxxx}
\StudentClass{BK21HTD}
\Major{Ky thuat Dien}
\Supervisor{TS. Nguyen Van B}
\ThesisMonth{06}
\ThesisYear{2026}
\end{lstlisting}

Sau khi chỉnh xong, biên dịch lại --- trang bìa, trang nhiệm vụ, tờ chấm sẽ \textbf{tự động cập nhật} toàn bộ thông tin.
